\documentclass[10pt,a4paper]{amsart}
\usepackage[margin=19mm]{geometry}
\usepackage{amsmath,amssymb,mathtools}
\usepackage[scr=rsfs,cal=euler]{mathalfa}
\usepackage{xfrac}
\usepackage[unicode,hidelinks,bookmarks=false]{hyperref}
\newcommand{\ie}{i.e.\ }
\newcommand{\cf}{cf.\ }
\newcommand{\resp}{respectively\ }
\newcommand{\Subsection}{\subsection}
\providecommand{\nobreakdash}{\nobreak\hbox{-}}
\setlength{\emergencystretch}{2em}
\multlinegap0pt
\usepackage{xcolor}
\usepackage{amssymb}
\usepackage{mathtools}
\usepackage{braket}
\usepackage{bm}
\usepackage{tikz}
\usepackage[shortlabels]{enumitem}
\usepackage{csquotes}
\newtheorem{lemma}{Lemma}
\newcommand{\F}{\mathcal{F}}
\newcommand{\Fh}[1]{\widehat{#1}}
\newcommand{\Hnxy}[3]{%
    \left\lVert #1\right\rVert_{H^{#2}_{#3}}%
}
\newcommand{\Lptxy}[3]{%
    \left\lVert #1\right\rVert_{L^{#2}_{#3}}%
}
\def\R{{\mathbb R}}
\def\T{{\mathbb T}}
\newcommand{\brs}[1]{%
    \left\{#1\right\}
}
\newcommand{\inp}[1]{%
    \left\langle #1\right\rangle
}
\newcommand{\norm}[1]{%
    \left\lVert #1\right\rVert%
}
\def\pt{\partial}
\newcommand{\sg}[1]{%
    e^{#1\partial_x\Delta}%
}
\newcommand{\xnorm}[3]{%
    \left\lVert #1\right\rVert_{X^{{#2},{#3}}}%
}
\title{On the modified Zakharov--Kuznetsov equation on the cylinder: A trilinear Bourgain-space estimate and its application to local well-posedness}
\author{Ali Mezher}
\date{Source: arXiv:2407.04850v3 (CC BY 4.0)}
\begin{document}
\maketitle

\noindent\emph{Source and licence.} On the modified Zakharov--Kuznetsov equation on the cylinder: A trilinear Bourgain-space estimate and its application to local well-posedness. Version arXiv:2407.04850v3 (CC BY 4.0), distributed by arXiv under the Creative Commons Attribution 4.0 International licence, http://creativecommons.org/licenses/by/4.0/, as recorded in the arXiv licence metadata for this exact version and shown on the abstract page https://arxiv.org/abs/2407.04850v3. Full licence notice: \texttt{licenses/arxiv-2407.04850-CC-BY-4.0.txt}.\par\smallskip

\noindent\textit{Editorial scope.} Counted unit: the article's lemma Lem:homogeneous-linear-estimate (source0.tex lines 4327--4335). The article's complete original proof of it is delivered with it (source0.tex lines 4338--4362, one \texttt{proof} environment of 25 lines). No mathematics is written, repaired or re-derived: the statement and the proof are the article's own lines, in the article's own notation, and no retained line is substituted or re-flowed.  No further source-local context is needed: every cross-reference of every retained line resolves inside the two delivered spans.  Nothing was omitted from any retained span, so the package carries no omission record.  No retained line contains a citation, so the wrapper carries no bibliography and no bibliography entry.  No retained line contains a figure environment, an image-inclusion command or drawing code, so no image is delivered and the package declares no asset dependency.  Editorial formatting only, in addition: an AMS \texttt{amsart} preamble that restates the article's own declarations and macros used by the retained lines, namely the environments \texttt{lemma} on separate counters, together with the standard packages those lines need.  The retained environments print in this document as Lemma 1 in order of appearance, whereas the article prints the same items with the numbering of its own full document.  The complete native source of the article as distributed in the arXiv e-print for this exact version is bundled unchanged as \texttt{sources/161886/source0.tex}.
\begin{lemma}[Homogeneous linear estimate] 
\label{Lem:homogeneous-linear-estimate}

Let $s,b\in\R$ and $u_0\in H^s(\R\times\T)$.
\begin{equation}
\label{lem:homogeneous-linear-estimate}
\xnorm{{\eta(t) e^{-t\pt_x\Delta}u_0}}{s}{b} \lesssim_b \Hnxy{u_0}{s}{xy}.
\end{equation}
\end{lemma}

\begin{proof}
    We start by calculating the space-time Fourier transform of $\eta\, \sg{-t} u_0$, 


    \begin{align*}
        \F {\left({\eta(t) \sg{-t} u_0}\right)}(\tau, \xi,q) &= \frac{1}{(2\pi)^3} \int e^{- i (t \tau + (x,y) \cdot (\xi,q))} \eta(t) e^{i t \omega(\xi,q)} u_0(x,y) dx dy dt\\
        & = \brs{ \frac{1}{2\pi} \int_\R e^{- i t (\tau - \omega(\xi,q))} \eta(t) dt}\\
        &\quad \times \brs{\frac{1}{(2 \pi)^2} \int_{\R \times \T} e^{- i (x,y)\cdot (\xi,q)} u_0(x,y) dx dy }\\
        &= \Fh{\eta} (\tau - \omega(\xi,q)) \Fh{u_0}(\xi,q).
    \end{align*}
Now we can calculate the norm,

\begin{align*}
    \xnorm{\eta(t) \sg{-t} {u_0}}{s}{b} &= \norm{\inp{|(\xi,q)|}^s \inp{\tau - \omega(\xi,q)}^b \Fh{\eta}(\tau -\omega(\xi,q)) \Fh{u_0}(\xi,q)}_{L_{\xi q}^2 L^2_\tau}\\
    &= \Lptxy{  \Lptxy{\inp{\tau -\omega(\xi,q)}^b \Fh{\eta}(\tau - \omega(\xi,q))}{2}{\tau} \inp{|(\xi,q)|}^s \Fh{u_0}(\xi,q)  }{2}{\xi q}\\
    &= \norm{\eta}_{H^b_t} \norm{\inp{|(\xi,q)|}^s \Fh{u_0}}_{L^2_{\xi q}}\\
    &= \norm{\eta}_{H^b_t} \norm{u_0}_{H^{s}_{xy}}\\
    & \lesssim_b \norm{u_0}_{H^{s}_{xy}}.
\end{align*}




    
\end{proof}

\end{document}
